\section{The subobject comes first}
\label{cyhall:sec:arrow}

An extension correspondence gives a product only after its
coefficient maps have been specified. For a quiver without relations,
these maps are ordinary pullback and complex-oriented integration.
A Calabi--Yau moduli space is generally singular. Its corresponding
operation uses a perverse sheaf of vanishing cycles and an orientation.
The one-arrow quiver fixes the order and grading before those
coefficients are introduced.

For the Hall constructions below, work over $\mathbb C$ and use
rational constructible coefficients.

For a finite quiver $Q$, write $Q_0$ for its vertices and $Q_1$
for its arrows. A representation consists of vector spaces $V_i$
and maps $V_a:V_i\to V_j$ for each arrow $a:i\to j$.
Its dimension vector is $\alpha=(\dim V_i)_{i\in Q_0}$.
Put
\[
 R_\alpha=\bigoplus_{a:i\to j}\operatorname{Hom}
             (\mathbb C^{\alpha_i},\mathbb C^{\alpha_j}),
 \qquad G_\alpha=\prod_iGL_{\alpha_i},\qquad
 \mathfrak M_\alpha=[R_\alpha/G_\alpha].
\]
Contraction of $R_\alpha$ to zero identifies
$H^\bullet(\mathfrak M_\alpha;\mathbb Q)$ with
$H^\bullet(BG_\alpha;\mathbb Q)$.
Its Chern roots have degree two.

An extension is written
\begin{equation}\label{cyhall:eq:extension}
 0\longrightarrow E\longrightarrow V\longrightarrow F
 \longrightarrow0.
\end{equation}
The ordered input is $(E,F)$: subobject, then quotient.
Let $\mathfrak E_{\alpha,\beta}$ classify
\eqref{cyhall:eq:extension} with dimensions $(\alpha,\beta)$.
There are maps
\begin{equation}\label{cyhall:eq:correspondence}
 \mathfrak M_\alpha\times\mathfrak M_\beta
 \xleftarrow{\ q\ }\mathfrak E_{\alpha,\beta}
 \xrightarrow{\ p\ }\mathfrak M_{\alpha+\beta},
 \qquad q(V,E)=(E,V/E),\quad p(V,E)=V.
\end{equation}
Over $\mathfrak M_{\alpha+\beta}$, the correspondence is the closed
incidence subscheme of the product of Grassmannian bundles
$\prod_i\operatorname{Gr}(\alpha_i,\alpha_i+\beta_i)$.
Hence $p$ is representable and proper.

Take $Q$ to have one arrow $1\to2$, and let $S_1,S_2$ be its
one-dimensional simple representations. Write $x,y$ for the Chern
classes at the two vertices, and $u_1,u_2$ for the constant classes
of their moduli stacks.
A representation of dimension $(1,1)$ is a scalar map $z:\mathbb C\to
\mathbb C$.

For subobject $S_1$ and quotient $S_2$, invariance forces $z=0$.
This is the zero section of the equivariant line of weight $y-x$.
For subobject $S_2$ and quotient $S_1$, every scalar $z$ is allowed.
Complex-oriented Gysin integration therefore gives
\begin{equation}\label{cyhall:eq:arrow-products}
                    u_1\star u_2=y-x,
                    \qquad u_2\star u_1=1
 \quad\text{in }H^\bullet(\mathfrak M_{(1,1)};\mathbb Q).
\end{equation}
The two products have ordinary degrees two and zero.
Their common target is the same stack, with different extension
conditions.

\subsection{Normal and tangent Euler classes}

For arbitrary $Q$, let $S_i$ and $T_i$ be the tautological subbundle
and quotient bundle on the product of Grassmannians above.
Preservation of the subspaces sets to zero every block
$S_i\to T_j$ along an arrow $i\to j$.
The normal bundle of this incidence condition and the tangent bundle
of the Grassmannians are respectively
\begin{equation}\label{cyhall:eq:normal-tangent}
 N_{\alpha,\beta}=\bigoplus_{a:i\to j}
                   \operatorname{Hom}(S_i,T_j),
 \qquad
 T_{\alpha,\beta}=\bigoplus_i\operatorname{Hom}(S_i,T_i).
\end{equation}
For classes $f$ and $g$, the Hall product is the Grassmannian
integral of $f(S)g(T)e(N_{\alpha,\beta})$.
Here $e$ is the top Chern class. Thus the construction is polynomial
before equivariant localization is used to compute it.

Define the quiver Euler form by
\begin{equation}\label{cyhall:eq:quiver-euler}
 \chi_Q(\alpha,\beta)=
 \sum_i\alpha_i\beta_i-\sum_{a:i\to j}\alpha_i\beta_j.
\end{equation}
The normal Gysin map raises degree by
$2\operatorname{rank}N_{\alpha,\beta}$, and integration lowers it by
$2\operatorname{rank}T_{\alpha,\beta}$.
Consequently
\begin{equation}\label{cyhall:eq:ordinary-degree}
 H^p(\mathfrak M_\alpha)\otimes H^q(\mathfrak M_\beta)
 \longrightarrow
 H^{p+q-2\chi_Q(\alpha,\beta)}(\mathfrak M_{\alpha+\beta}).
\end{equation}
This is the ordinary cohomological Hall construction
\cite{cyhall:KS}*{\S\S2.4--2.5}.

For its localized expression, partition the roots at each vertex
into $I_i\sqcup J_i$, with $|I_i|=\alpha_i$ and $|J_i|=\beta_i$.
Each such partition contributes
\begin{equation}\label{cyhall:eq:shuffle}
 f(x_I)g(x_J)
 \frac{\displaystyle\prod_{a:i\to j}
        \prod_{r\in I_i,\ s\in J_j}(x_{j,s}-x_{i,r})}
      {\displaystyle\prod_i
        \prod_{r\in I_i,\ s\in J_i}(x_{i,s}-x_{i,r})}.
\end{equation}
Summing over all partitions gives $f\star g$.
The numerator records precisely the forbidden arrow blocks.
The denominator records the tangent directions of the flag.

For three inputs, both bracketings sum over the same ordered
three-block partitions. Every pair of distinct blocks contributes
its factor once. This proves associativity of the displayed
formula and agrees with composition of three-step flags.

\subsection{Perverse normalization}

For a cochain complex $K$ we use $K[s]^n=K^{n+s}$.

The smooth stack $\mathfrak M_\alpha$ has dimension
\[
 d_\alpha=\dim R_\alpha-\dim G_\alpha
          =-\chi_Q(\alpha,\alpha).
\]
Its constant perverse sheaf is
$\mathbb Q_{\mathfrak M_\alpha}[d_\alpha]$.
Set
\[
 \mathsf H^p_\alpha=
 H^p(\mathfrak M_\alpha,\mathbb Q[d_\alpha])
 =H^{p+d_\alpha}(\mathfrak M_\alpha;\mathbb Q).
\]

Changing to this grading changes the degree in
\eqref{cyhall:eq:ordinary-degree} to
\begin{equation}\label{cyhall:eq:perverse-degree}
 \begin{split}
 \nu_Q(\alpha,\beta)
 &=-2\chi_Q(\alpha,\beta)+d_\alpha+d_\beta-d_{\alpha+\beta}\\
 &=\chi_Q(\beta,\alpha)-\chi_Q(\alpha,\beta).
 \end{split}
\end{equation}
In the one-arrow example, all three stack dimensions are $-1$.
Both constant input classes have perverse degree one.
The outputs $y-x$ and $1$ have perverse degrees three and one,
respectively. Their product degrees are $+1$ and $-1$.
This also fixes the shift when the coefficient sheaf is replaced
by vanishing cycles.

\section{Flags and derived composition}
\label{cyhall:sec:flags}

For three objects, composing extensions gives
$0\subset E_1\subset E_2\subset E_3$.
One bracketing first forms $E_2$, the other $E_3/E_1$.
Both retain the same three successive quotients.
In a derived category the quotient is a cofiber, so its nullhomotopy
is part of the flag.

Let $\mathcal D$ be a small stable $\infty$-category.
Thus it has a zero object and finite limits and colimits, and a
square is a pullback precisely when it is a pushout.

For $n\geq0$, let $T_n$ be the poset of pairs $(i,j)$ with
$0\leq i\leq j\leq n$, ordered coordinatewise.
Let $S_n(\mathcal D)$ be the space of coherent $T_n$-diagrams
\[
               X_{ij},\qquad 0\leq i\leq j\leq n,
\]
with $X_{ii}=0$, such that, for $i\leq j\leq k$, the square
\[
 \begin{tikzcd}[column sep=large]
 X_{ij}\ar[r]\ar[d]&X_{ik}\ar[d]\\
 0\ar[r]&X_{jk}
 \end{tikzcd}
\]
is a pushout and a pullback.
Equivalences of diagrams give the paths in this space.
Restriction to a subset of the indices gives its face maps.
Repetition of an index inserts a zero quotient and gives its
degeneracy maps. This is the Waldhausen $S$-construction.

\begin{proposition}[Gluing flags]
\label{cyhall:prop:two-segal}
Triangulate the polygon with vertices $0,\ldots,n$.
Restriction identifies $S_n(\mathcal D)$ with the homotopy limit
of its triangle spaces, glued along common edges.
In particular, both extension composites are equivalent
to $S_3(\mathcal D)$. Inserting a zero quotient gives the unit.
\end{proposition}
\begin{proof}
For the diagonal $02$ of the quadrilateral $0123$, the two
triangles specify $X_{01}\to X_{02}\to X_{03}$ together with
the corresponding cofibers. The remaining object is
$X_{13}=\operatorname{cofib}(X_{01}\to X_{03})$.
Pasting pushout squares gives the required map
$X_{12}\to X_{13}$ and its cofiber $X_{23}$.
The space of these completions is contractible by the universal
property of a colimit.

For the diagonal $13$, the triangles specify
$X_{01}\to X_{03}\to X_{13}$ and
$X_{12}\to X_{13}\to X_{23}$.
Recover $X_{02}$ as the pullback of $X_{03}\to X_{13}$ along
$X_{12}\to X_{13}$. Stability makes this square a pushout as well,
so it recovers the same flag.
Again the completion is unique up to a contractible choice.

Gluing triangles successively and applying these universal
properties proves the assertion for a polygon.
The coherent version is the unital $2$-Segal theorem for the
$S$-construction, \cite{cyhall:DK}*{Theorem~7.3.3}.
A stable category has the exact structure in which all maps
are admissible.
A zero cofiber makes its map an equivalence, which proves the
unit assertion.
\end{proof}

The proposition concerns spaces of diagrams, including their
automorphisms and higher homotopies. Applied to stable categories
of families over a derived test algebra, it gives equivalences
of prestacks. Its fiber products are derived.

For example, put $A=\mathbb C[t]$ and let both maps
$\operatorname{Spec}\mathbb C\to\operatorname{Spec}A$ be the origin.
Resolving $\mathbb C$ by
$[A\xrightarrow{t}A]$ in degrees $-1,0$ gives
\begin{equation}\label{cyhall:eq:self-intersection}
 \mathbb C\otimes_A^{\mathrm L}\mathbb C
       \simeq\Lambda_{\mathbb C}(\eta),
       \qquad |\eta|=-1,\qquad d\eta=0.
\end{equation}
The ordinary fiber product retains only degree zero.
The missing class belongs to the composed correspondence itself,
before a coefficient theory is applied.

\section{The Calabi--Yau coefficient map}
\label{cyhall:sec:coefficients}

Let $\mathcal C$ be a dg category of finite type over $\mathbb C$,
equipped with a left Calabi--Yau structure of dimension three.
Here finite type means compactness in the Morita category, and
a Calabi--Yau structure is a nondegenerate negative cyclic class.
Write $\mathcal M_{\mathcal C}$ for its derived moduli stack of
pseudo-perfect right modules: an $R$-family is a dg functor
$\mathcal C^{\mathrm{op}}\to\operatorname{Perf}(R)$.
It carries the resulting $(-1)$-shifted symplectic form.

Choose an open substack $\mathcal M\subset\mathcal M_{\mathcal C}$
closed under direct sums and containing the zero module.
Put $\Theta=[\mathbb A^1/\mathbb G_m]$, where
$\mathbb G_m$ acts by scaling, and set
\[
 \operatorname{Filt}(\mathcal M)=\operatorname{Map}(\Theta,\mathcal M),
 \qquad
 \operatorname{Grad}(\mathcal M)=\operatorname{Map}(B\mathbb G_m,\mathcal M).
\]
Restriction at $1$ gives evaluation, and restriction at zero gives
the associated graded object.
These definitions include the automorphisms of a filtration.

We impose the following geometric hypotheses:
\begin{enumerate}
 \item The classical truncation $\mathcal M^{\mathrm{cl}}$ is a
 quasi-separated Artin stack with affine stabilizers.
 \item The evaluation map
 $\operatorname{Filt}(\mathcal M^{\mathrm{cl}})\to
 \mathcal M^{\mathrm{cl}}$ is proper on each connected component
 of its source.
 \item The zero-module map
 $\epsilon:\operatorname{Spec}\mathbb C\to\mathcal M^{\mathrm{cl}}$
 is a closed immersion.
\end{enumerate}
The second condition is called \emph{$\Theta$-reductivity} here.

Put $\Gamma=\pi_0(\mathcal M)$, with addition induced by direct sum,
and write $\mathcal M_\alpha$ for a component.
Map $\mathcal M_\alpha\times\mathcal M_\beta$ to
$\operatorname{Grad}(\mathcal M)$ by giving $E$ weight one and
$F$ weight zero. Define
\[
 \mathcal E_{\alpha,\beta}=
 \operatorname{Filt}(\mathcal M)
 \mathop{\times}^{\mathrm R}_{\operatorname{Grad}(\mathcal M)}
 (\mathcal M_\alpha\times\mathcal M_\beta).
\]
This is the correspondence of cofiber sequences $E\to V\to F$
with those weights. For modules in a heart it is the Rees
convention for $0\subset E\subset V$.
Projection and evaluation give
\[
 \mathcal M_\alpha\times\mathcal M_\beta
 \xleftarrow{\ q\ }\mathcal E_{\alpha,\beta}
 \xrightarrow{\ p\ }\mathcal M_{\alpha+\beta},
\]
with the same subobject-first order as
\eqref{cyhall:eq:correspondence}.

\subsection{Strong orientation}

An orientation $o$ consists of a $\mathbb Z/2$-graded line $L_o$
on $\mathcal M^{\mathrm{red}}$ and a squaring isomorphism
\[
 L_o^{\otimes2}\simeq
 \det(\mathbb L_{\mathcal M})|_{\mathcal M^{\mathrm{red}}},
\]
where the determinant has its virtual-rank parity.
This is \cite{cyhall:KPS}*{Definition~3.15}.
On a critical chart the orientation determines the
sign local system needed to glue perverse vanishing cycles.
The resulting Donaldson--Thomas sheaf will be denoted $\mathcal P_o$.
It is a rational perverse sheaf on $\mathcal M^{\mathrm{cl}}$.
We use the vanishing-cycle normalization for which the zero
function on a smooth space has coefficient $\mathbb Q[\dim]$.

The direct-sum map $\Phi:\mathcal M\times\mathcal M\to\mathcal M$
has an oriented pullback operation $\Phi^\star$ supplied by its
filtered Lagrangian correspondence. This operation includes the
determinant of the moving tangent directions in addition to
ordinary pullback.
\begin{definition}\label{cyhall:def:strong-orientation}
A \emph{strong orientation datum} consists of $o$, an isomorphism
$\lambda:\Phi^\star o\simeq o\boxtimes o$, and a trivialization
$\epsilon^\star o\simeq o_{\mathrm{triv}}$ at the zero module.
The two isomorphisms obtained by applying $\lambda$ to three
summands must agree under the canonical comparison of the
filtered Lagrangian correspondences. The left and right unit
comparisons must agree with the given trivialization.
\end{definition}
This is \cite{cyhall:KPS}*{Definition~8.4}.
It specifies the compatibility isomorphisms themselves.
Existence of a square root at each object does not give these
isomorphisms or their associativity.

\subsection{The index and the product}

At $E\oplus F$, decompose the tangent complex into
$\mathbb G_m$-weights. Define the index
\begin{equation}\label{cyhall:eq:index}
 I(\alpha,\beta)=
 \operatorname{rank}\bigl(\mathbb T_{\mathcal M,E\oplus F}^{>0}\bigr),
 \qquad
 \operatorname{rank}K=\sum_j(-1)^j\dim H^j(K).
\end{equation}
This index is locally constant on the graded component.
Equivalently, take the positive-minus-negative weight dimension
in $H^0$ of the classical tangent complex.
Subtract the same difference in $H^{-1}$
\cite{cyhall:KPS}*{Definition~6.6 and (6.4)}.

At objects with perfect derived Hom complexes, put
$\chi(E,F)=\sum_j(-1)^j\dim\operatorname{Ext}^j(E,F)$.
The positive-weight tangent is $\operatorname{RHom}(F,E)[1]$.
Three-dimensional Calabi--Yau duality therefore gives
\begin{equation}\label{cyhall:eq:index-euler}
            I(\alpha,\beta)=-\chi(F,E)=\chi(E,F).
\end{equation}
The shift by one in the tangent is essential to this sign.
It also shows that $I$ is additive in each charge whenever the
charges are taken in the Grothendieck group.

On connected components it satisfies the identity
\begin{equation}\label{cyhall:eq:cocycle}
 I(\alpha,\beta)+I(\alpha+\beta,\gamma)
 =I(\beta,\gamma)+I(\alpha,\beta+\gamma).
\end{equation}
This identity also follows directly by comparing the moving
directions of a three-step filtration.

\begin{theorem}[Kinjo--Park--Safronov]
\label{cyhall:thm:coefficient}
Under the stated geometric hypotheses and a strong orientation
datum, write $\mathcal P_\alpha=\mathcal P_o|_{\mathcal M_\alpha}$.
There is a natural coefficient isomorphism
\begin{equation}\label{cyhall:eq:coefficient}
 \zeta_{\alpha,\beta}:
 (\mathcal P_\alpha\boxtimes\mathcal P_\beta)[I(\alpha,\beta)]
 \xrightarrow{\ \sim\ }\mathrm Rq_*p^!\mathcal P_{\alpha+\beta}.
\end{equation}
Its two composites on a three-step filtered correspondence agree,
and its zero-object restrictions are the identities.
Consequently the spaces
$\mathsf H_\alpha^p=H^p(\mathcal M_\alpha,\mathcal P_\alpha)$
have a unital associative Hall multiplication
\begin{equation}\label{cyhall:eq:critical-product}
 \mathsf H_\alpha^p\otimes\mathsf H_\beta^q
 \longrightarrow
 \mathsf H_{\alpha+\beta}^{p+q-I(\alpha,\beta)}.
\end{equation}
\end{theorem}

The coefficient theorem is
\cite{cyhall:KPS}*{Theorem~8.6 and Proposition~8.29},
with the two graded factors ordered as above.
Its construction uses hyperbolic localization of oriented vanishing
cycles. The following adjunction calculation specifies how it gives
the product and fixes the target shift.

\begin{proof}[Construction of the product]
Write $C_\alpha=\mathrm R\Gamma(\mathcal M_\alpha,\mathcal P_\alpha)$.
First take the external product of sections. Then use
\eqref{cyhall:eq:coefficient}:
\[
 \begin{split}
 (C_\alpha\otimes C_\beta)[I(\alpha,\beta)]
 &\longrightarrow
 \mathrm R\Gamma(\mathcal M_\alpha\times\mathcal M_\beta,
       (\mathcal P_\alpha\boxtimes\mathcal P_\beta)[I])\\
 &\xrightarrow{\ \zeta_{\alpha,\beta}\ }
 \mathrm R\Gamma(\mathcal E_{\alpha,\beta},p^!\mathcal P_{\alpha+\beta}).
 \end{split}
\]
The canonical external product suffices, without a K\"unneth
isomorphism for an unrestricted infinite stack.

Properness of $p$ identifies $\mathrm Rp_*$ with $\mathrm Rp_!$.
The counit $\mathrm Rp_!p^!\mathcal P\to\mathcal P$ gives
\begin{equation}\label{cyhall:eq:chain-product}
 (C_\alpha\otimes C_\beta)[I(\alpha,\beta)]
 \longrightarrow C_{\alpha+\beta},
 \quad\text{equivalently}\quad
 C_\alpha\otimes C_\beta\longrightarrow
 C_{\alpha+\beta}[-I(\alpha,\beta)].
\end{equation}
Taking cohomology gives \eqref{cyhall:eq:critical-product}.

For three inputs, identify both composed correspondences with the
three-step filtered correspondence. Proper base change and
composition of counits identify the two final integrations.
The two coefficient maps agree by the associativity part of the
coefficient theorem. Equation~\eqref{cyhall:eq:cocycle} makes their
shifts equal, so the induced cohomological products agree with
the usual Koszul signs for shifted complexes.
The zero module, its orientation trivialization and the unitality
of $\zeta$ give the unit, as in
\cite{cyhall:KPS}*{Corollary~8.8}.
\end{proof}

The source shift in \eqref{cyhall:eq:coefficient} determines the
minus sign in \eqref{cyhall:eq:chain-product}.
For a quiver with potential, the comparison with the critical
quiver construction is
\cite{cyhall:KPS}*{Proposition~8.37}.
With subobject first, its index is
$I(\alpha,\beta)=\chi_Q(\alpha,\beta)-\chi_Q(\beta,\alpha)$.
Thus \eqref{cyhall:eq:critical-product} recovers
\eqref{cyhall:eq:perverse-degree}.
Reversing the order of the inputs reverses this antisymmetric
index. The one-arrow products \eqref{cyhall:eq:arrow-products}
provide the elementary normalization.

\section{Coherent coefficients and local flag calculations}
\label{cyhall:sec:coherent}

The $2$-Segal statement identifies spaces of composable flags.
The coefficient theorem supplies the maps needed for cohomological
Hall multiplication. To obtain an algebra of complexes with all
associativity homotopies, retain coefficient maps and their
coherences on every higher flag.

Here is a sufficient formulation. Take a class of flag
correspondences closed under the derived fiber products used in
composition. On $\Gamma$-graded complexes define the tensor product
\[
 (V\mathbin{\odot}W)_\delta=
 \bigoplus_{\alpha+\beta=\delta}
       (V_\alpha\otimes W_\beta)[I(\alpha,\beta)].
\]
Equation~\eqref{cyhall:eq:cocycle} gives its associator, using the
standard signed identifications for shifts and tensor products.
Assume a monoidal functor from the flag
$\infty$-category of correspondences to this category of complexes.

Require that it sends the extension span to
\eqref{cyhall:eq:chain-product} and the zero span to its unit.
Such a functor includes coherent base-change maps, composition
maps and unit maps. Equality of the induced cohomology maps is
weaker data.

\begin{proposition}[Conditional coherent realization]
\label{cyhall:prop:coherent-realization}
A coefficient functor with these properties sends the unital
$2$-Segal flag object to a unital $E_1$-algebra in the specified
category of charge-graded complexes.
\end{proposition}
\begin{proof}
Each triangulation gives an iterated multiplication span.
Proposition~\ref{cyhall:prop:two-segal} identifies it with the span
of the full flag. These identifications are equivalences of
coherent diagrams, not choices of isomorphism classes.

Applying the assumed functor carries them and their compositions
to the associativity coherences. Inserting a zero quotient gives
the unit coherences. The cocycle identity
\eqref{cyhall:eq:cocycle} identifies their shifts.
These are the operations and coherences of a unital $E_1$-algebra.
\end{proof}

On compact threefolds, an orientation compatible with direct sums
must still satisfy the strong associativity condition before the
coefficient theorem applies. A coherent realization must further
retain the required higher coefficient comparisons.
To produce holomorphic fields, one must in addition construct the
coefficient system on ordered collision strata and compare its
operations with the Hall maps. These are distinct constructions.

A local flag family can test those comparisons. Its coherent
direct image records functions and derived specialization. Its
constructible direct image records supported cohomology maps.
The Hall integration map is determined by
\eqref{cyhall:eq:coefficient} and its proper counit.
Neither a scheme-cycle multiplicity nor a Frobenius functional
on the coherent algebra identifies that trace without a
comparison of these coefficient operations.
